\Article{9}
Wie das vorige Verfahren auf der Zerlegung des Raumbereichs $S$ in prismenförmige Elemente beruht, so geht das zweite von der Zerlegung desselben Raumbereichs in pyramidenförmige Elemente aus. Man denke sich eine Kugeloberfläche vom Radius $=1$ um den Mittelpunkt $\mu$ beschrieben und in unendlich kleine Elemente zerlegt. Zu einem Punkt eines solchen Elements $\mathrm{d}\Pi$ ziehe man eine Gerade vom Punkt $\mu$ aus. Sie schneide die Oberfläche des Raumbereichs $S$ der Reihe nach in den Punkten $P',P'',P'''$ usw. Die Abstände dieser Punkte von $\mu$ seien mit $R',R'',R'''$ usw. bezeichnet. Die von $\mu$ zu allen Punkten des Randes des Elements $\mathrm{d}\Pi$ gezogenen Geraden bilden einen pyramidenförmigen Raumbereich und schneiden bei den Punkten $P',P'',P'''$ usw. Elemente aus der Oberfläche des Raumbereichs $S$ heraus, die wir mit $\mathrm{d}T',\mathrm{d}T'',\mathrm{d}T'''$ usw. bezeichnen. Schließlich sei $Q'$ der Winkel zwischen der Geraden $P'\mu$ und der nach außen gerichteten Normalen auf dem Element $\mathrm{d}T'$. Ebenso seien $Q'',Q'''$ usw. die Neigungen der entsprechenden Normalen bei den Punkten $P'',P'''$ usw. gegen die nach $\mu$ gezogene Gerade. Dann gilt
\[
\mathrm{d}\Pi=\pm\frac{\mathrm{d}T'.\cos Q'}{R'R'}=\mp\frac{\mathrm{d}T''.\cos Q''}{R''R''}=\pm\frac{\mathrm{d}T'''.\cos Q'''}{R'''R'''}\ \text{etc.}
\]
wobei die oberen oder unteren Vorzeichen gelten, je nachdem, ob $\mu$ außerhalb oder innerhalb des Raumbereichs $S$ liegt. Der Fall, dass $\mu$ auf der Oberfläche des Raumbereichs $S$ selbst liegt, ist zum ersten oder zum zweiten Fall zu rechnen, je nachdem, ob die Gerade $\mu P'$ außerhalb oder innerhalb des Raumbereichs $S$ verläuft.

Ferner ist klar, dass für alle Teile des Raumbereichs $S$ innerhalb jenes pyramidenförmigen Raumbereichs der Winkel $q$ konstant ist. Ähnlich wie in Artikel 7 folgern wir daher, dass, wenn für das unbestimmte Integral
\[
\int\psi r.\mathrm{d}r=-\theta r
\]
mit beliebig gewählter Integrationskonstante gesetzt wird, das Integral
\[
\int\frac{\mathrm{d}S.\cos q.\psi(\mu,\mathrm{d}S)}{(\mu,\mathrm{d}S)^2}
\]
über alle Teile des Raumbereichs $S$ innerhalb jenes pyramidenförmigen Raumbereichs im ersten Fall
\SourceLeaf{48}
\[
=\cos q.\left(\frac{\mathrm{d}T'.\cos Q'.\theta R'}{R'R'}+\frac{\mathrm{d}T''.\cos Q''.\theta R''}{R''R''}+\frac{\mathrm{d}T'''.\cos Q'''.\theta R'''}{R'''R'''}+\ \text{etc.}\right)
\]
wird, während im zweiten Fall zu derselben Formel noch der Term $\mathrm{d}\Pi.\cos q.\theta0$ hinzukommt.

Summieren wir nun über alle Elemente der Kugeloberfläche, so ergibt sich für das vollständige Integral
\[
\int\frac{\mathrm{d}S.\cos q.\psi(\mu,\mathrm{d}S)}{(\mu,\mathrm{d}S)^2}
\]

I. im Fall, dass der Punkt $\mu$ außerhalb des Raumbereichs $S$ liegt,
\[
=\int\frac{\mathrm{d}T.\cos q.\cos Q.\theta R}{RR}
\]
wobei $\mathrm{d}T$ allgemein die Elemente der Oberfläche des Raumbereichs $S$ bezeichnet und $Q,R$ für diese Elemente dieselben Größen bezeichnen wie zuvor die mit Strichen versehenen Buchstaben für bestimmte Elemente. Schließlich bezeichnet $q$ die Neigung der vom Punkt $\mu$ zum Element $\mathrm{d}T$ gezogenen Geraden gegen unsere feste Gerade.

II. Im Fall, dass der Punkt $\mu$ innerhalb des Raumbereichs $S$ liegt, ist der Term
\[
\theta0.\int\mathrm{d}\Pi.\cos q
\]
hinzuzufügen. Dabei ist $q$ die Neigung der von $\mu$ zu $\mathrm{d}\Pi$ gezogenen Geraden gegen die feste Gerade, und die Integration ist über die gesamte Kugeloberfläche zu erstrecken. Man sieht jedoch leicht, dass dieses Integral über die Halbkugel, für die $q$ spitz ist, den Wert $=+\pi$ und über die andere Halbkugel den Wert $=-\pi$ hat. Daher verschwindet das vollständige Integral, und für diesen zweiten Fall gilt genau dieselbe Formel, die wir für den ersten angegeben haben. Anders verhält es sich im dritten Fall.

III. Wenn der Punkt $\mu$ auf der Oberfläche des Raumbereichs $S$ selbst liegt. Auch hier ist der Term
\[
\theta0.\int\mathrm{d}\Pi.\cos q
\]
hinzuzufügen, doch wird nur über diejenigen Teile der Kugeloberfläche integriert, für die das Anfangsstück der von $\mu$ zu $\mathrm{d}\Pi$ gezogenen Geraden innerhalb des Raumbereichs $S$ liegt. Sofern die Oberfläche des Raumbereichs $S$ am Punkt $\mu$ weder eine Spitze noch eine Kante aufweist, sind dies die Teile, für die diese Gerade mit der nach außen gerichteten Normalen auf der Oberfläche des Raumbereichs $S$ am Punkt $\mu$ einen stumpfen Winkel bildet. Es bleibt also, das in diesem Sinn verstandene Integral zu bestimmen.

\SourceLeaf{49}Diese Normale und die feste Gerade mögen die Kugeloberfläche jeweils in den Punkten $G,H$ schneiden. Es sei der Bogen $GH=k$ und der Bogen zwischen $G$ und einem beliebigen Punkt der Kugeloberfläche $=v$. Schließlich sei $w$ der sphärische Winkel zwischen den Bögen $k,v$. Damit gilt $\cos q=\cos k.\cos v+\sin k.\sin v.\cos w$, und für $\mathrm{d}\Pi$ ist das Element $\sin v.\mathrm{d}v.\mathrm{d}w$ einzusetzen. Das Integral $\int\mathrm{d}\Pi.\cos q$
\[
=\int\int(\cos k.\cos v+\sin k.\sin v.\cos w)\sin v.\mathrm{d}v.\mathrm{d}w
\]
ist von $w=0$ bis $w=360^0$ und von $v=90^0$ bis $v=180^0$ zu erstrecken. So liefert die erste Integration
\[
\int2\pi\cos k.\cos v.\sin v.\mathrm{d}v
\]
und die zweite dann $-\pi\cos k$.

Für unseren Zweck kommt dieser dritte Fall nur insofern in Betracht, als die Oberflächen der Raumbereiche $s,S$ einen gewissen endlichen Teil gemeinsam haben. Liegt der Punkt $\mu$ auf diesem Teil, so gilt entweder $k=0$ oder $=180^0$. Das Integral $\int\mathrm{d}\Pi.\cos q$ ist damit entweder $=-\pi$ oder $=+\pi$, je nachdem, ob beim Punkt $\mu$ die Raumbereiche $s,S$ auf derselben Seite oder auf entgegengesetzten Seiten der beiden Oberflächen gemeinsamen Tangentialebene liegen.

Wenden wir dies auf unser ursprüngliches Integral $\int\int\mathrm{d}s.\mathrm{d}S.\varphi(\mathrm{d}s,\mathrm{d}S)$ an, so ergibt sich sein Wert

I. wenn die Oberflächen der Raumbereiche $s,S$ keinen endlichen Teil gemeinsam haben,
\[
=4\pi\sigma\psi0+\int\int\frac{\mathrm{d}t.\mathrm{d}T.\cos q.\cos Q.\theta(\mathrm{d}t,\mathrm{d}T)}{(\mathrm{d}t,\mathrm{d}T)^2}
\]

II. wenn die Oberflächen der Raumbereiche $s,S$ einen endlichen Teil $=\tau$ gemeinsam haben,
\[
=4\pi\sigma\psi0\mp\pi\tau\theta0+\int\int\frac{\mathrm{d}t.\mathrm{d}T.\cos q.\cos Q.\theta(\mathrm{d}t,\mathrm{d}T)}{(\mathrm{d}t,\mathrm{d}T)^2}
\]
wobei das obere oder untere Vorzeichen gilt, je nachdem, ob die Raumbereiche $s,S$ auf derselben Seite oder auf entgegengesetzten Seiten der gemeinsamen Oberfläche $\tau$ liegen.

III. Wenn die Oberflächen der Raumbereiche $s,S$ mehrere getrennte endliche Teile gemeinsam haben, sei $\tau$ die Summe der Teile, an die die Raumbereiche $s,S$ von derselben Seite her angrenzen, und $\tau'$ die Summe der Teile, an die diese Raumbereiche von entgegengesetzten Seiten her angrenzen. Unser Integral wird dann
\SourceLeaf{50}
\[
=4\pi\sigma\psi0+\pi(\tau'-\tau)\theta0+\int\int\frac{\mathrm{d}t.\mathrm{d}T.\cos q.\cos Q.\theta(\mathrm{d}t,\mathrm{d}T)}{(\mathrm{d}t,\mathrm{d}T)^2}
\]
Diese dritte Formel lässt sich als Zusammenfassung aller Fälle ansehen. Das doppelte Integral ist über sämtliche Elemente beider Oberflächen zu erstrecken. Dabei bezeichnen $q,Q$ die Winkel, die die Verbindungslinie der beiden Elemente $\mathrm{d}t,\mathrm{d}T$ mit den nach außen gerichteten Normalen auf diesen Elementen bildet. Für den ersten Winkel wird die Richtung dieser Linie von $\mathrm{d}t$ nach $\mathrm{d}T$ genommen, für den zweiten von $\mathrm{d}T$ nach $\mathrm{d}t$.

\Article{10}
Die beiden in den Artikeln 8 und 9 entwickelten Umformungen des Integrals $\int\mathrm{d}s.\mathrm{d}S.\varphi(\mathrm{d}s,\mathrm{d}S)$ zeichnen sich durch beinahe dieselbe Eleganz aus, doch eignet sich die zweite besser für unseren Zweck. Die allgemeine Aufgabe lässt sich nur weiter reduzieren, wenn wir bestimmte Annahmen über die Raumbereiche $s,S$ oder über die Funktion $\varphi$ treffen. Da die Funktion $\varphi$ aus der Funktion $f$ hervorgeht, legen wir der weiteren Untersuchung dieselbe Hypothese zugrunde, von der der bedeutende \textsc{Laplace} ausgegangen ist, nämlich dass molekulare Anziehungskräfte nur bei unmerklichen Entfernungen merklich sind. Da dieser Ausdruck etwas unbestimmt bleibt, solange keine Einheit angegeben ist, bemerken wir zunächst, dass die Anziehungskraft $fr$, als Funktion der Entfernung $r$ ausgedrückt, erst mit einer Masse multipliziert werden muss, damit sie dieselbe Dimension wie die Schwerkraft $g$ hat. Die Annahme besagt dann Folgendes. Bezeichnet $M$ eine Masse, wie sie in Experimenten vorkommen kann und gegenüber der ganzen Erde als verschwindend gelten darf, so muss $Mfr$ gegenüber der Schwerkraft stets unmerklich bleiben, solange $r$ einen noch so kleinen, unseren Messungen zugänglichen Wert hat. Nichts schließt jedoch aus, dass der Wert von $Mfr$ bei unmerklichen Entfernungen merklich wird und bei abnehmendem $r$ sogar alle Grenzen überschreitet. Wir stellen gewiss nicht ohne Bewunderung fest, welch gewichtige Ergebnisse von ganz eigenem mathematischem Charakter sich aus dieser einen Hypothese gewinnen lassen, obwohl das Gesetz der Funktion $fr$ im Übrigen als völlig unbekannt gilt. Zwar können sie unter diesen Umständen keine absolute mathematische Genauigkeit beanspruchen, doch besitzen sie mit Sicherheit eine solche Genauigkeit, dass kein Experiment irgendeine Abweichung von der absoluten Wahrheit feststellen kann. Denn sobald es gelänge, eine solche Abweichung einer Messung zu unterwerfen, wäre die Annahme selbst hinfällig.

\SourceLeaf{51}\Article{11}
Wir dürfen annehmen, dass die Funktion $fr$, ebenso wie die Funktion $Fr$, die Anziehung unter Auslassung desjenigen Anteils bezeichnet, der zu $rr$ umgekehrt proportional ist und zur Erklärung der astronomischen Erscheinungen dient. Denn dieser Anteil kann, wie auch immer die Flüssigkeit und das Gefäß geformt sind, an jedem Punkt nur eine unmerkliche Änderung der Schwerkraft bewirken. Wächst also $r$ von einem merklichen Wert ins Unendliche, so ist $fr$ an sich unmerklich und nimmt außerdem schneller ab als $\frac{1}{rr}$. Daraus folgt leicht, dass auch das Integral $\int fr.\mathrm{d}r$, von einem beliebigen merklichen Wert bis ins Unendliche erstreckt, unmerklich ist. Deshalb nehmen wir an, dass die Integrationskonstante in $\int fr.\mathrm{d}r=-\varphi r$ so gewählt ist, dass $\varphi\infty=0$ gilt, oder dass $\varphi r$ der Wert des Integrals $\int fx.\mathrm{d}x$ von $x=r$ bis $x=\infty$ selbst ist. Auf diese Weise bezeichnet $\varphi r$ für jede Entfernung $r$ eine positive Größe, die unmerklich ist, solange $r$ merklich ist. Bei einem unmerklichen Wert von $r$ kann sie dagegen merklich werden und bei stetig abnehmender Entfernung $r$ sogar alle Grenzen überschreiten. In der üblichen Ausdrucksweise steht also nichts im Wege, dass $\varphi0=\infty$ gilt.

